\section{Codimension levels and modular dimension levels}
\label{sec:levels}

The manuscript uses both a filtration by normal-crossing codimension and
a filtration by modular dimension. They have different index posets and
connectivity estimates. A map is \(r\)-connected here when its homotopy
fibres are \((r-1)\)-connected; it induces isomorphisms on \(\pi_i\) for
\(i<r\) and an epimorphism on \(\pi_r\). A \(k\)-type equivalence is an
equivalence after \(\tau_{\leq k}\).

\subsection{The diagram (246), its coherence, and the sharp cutoff}

Let \(L=\partial X^{\log}\) be the boundary of the real oriented blow-up
of an NC pair. Define \(L_{\leq r}\) as the inverse image of the locus
where at most \(r\) local branches meet. Thus the truncation is an actual
subspace, with every incidence flag retained. For a proper pair a rounded
collar identifies \(L\) with a deleted neighbourhood of the divisor in
the open part. For a nonproper pair we use the resolved boundary itself;
an arbitrary end of the open space need not be such a collar.

\begin{proposition}[The two-level reading of (246)]\label{P5a1}
\status{Refuted.}
The bicategorical colimit in groupoids of the depth-one/depth-two
incidence span displayed in (246) need not be \(\Pi_1(L)\).
For the three coordinate hyperplanes in \(\C^3\), that colimit has vertex
group \(\Z^4\), whereas \(\pi_1(L)=\Z^3\). The depth-three object and its
flag triangles remove the extra generator.
\end{proposition}

\begin{proof}
Polar coordinates give
\[
 L=(S^1)^3\times
 \{(r_1,r_2,r_3)\in\R_{\geq0}^3:\min r_i=0\}.
\]
The radial factor contracts to the origin, so \(L\simeq(S^1)^3\).
Depth at most two deletes that origin. Dividing by the positive radial
coordinate identifies the remaining radial factor with the boundary
of a two-simplex. Hence \(L_{\leq2}\simeq(S^1)^3\times S^1\).

This also audits the actual two-level incidence diagram. Its indices
are the three singleton and three two-element subsets of
\(\{1,2,3\}\), with six inclusions, forming a hexagon. Each total-link
space is \((S^1)^3\) up to homotopy and every map induces the identity
on its three coordinate meridians. The bicategorical groupoid colimit
is the fundamental groupoid of the homotopy colimit of these classifying
spaces. The hexagon therefore supplies an additional infinite cyclic
generator. A spanning-tree presentation has one index generator commuting
with the three meridians.

Adjoining \(\{1,2,3\}\) and all composable flags gives the barycentric
subdivision of a filled two-simplex. The six triangular relations
eliminate the six non-tree edges of its enlarged graph.
The program \texttt{levels\_colimit.py} computes both presentations by
single-occurrence Tietze eliminations. Its outputs are
\((V,E,F)=(6,6,0)\) and \((7,12,6)\), with one and zero surviving
index generators respectively.

The strict colimit of the corresponding one-object groups would already
be \(\Z^3\). That strict colimit discards the path and natural-isomorphism
data required for descent and does not interpret the bicategorical
colimit in (246). The fully coherent flag diagram includes the
depth-three triangles and gives the correct answer.
\end{proof}

The marginal reference to \(2\)-fl\`eches on folio 76, scan 152,
supports recording the natural isomorphisms on triple flags. It does
not make an unwritten triangular relation automatic. The neighbouring
purity sentence cannot justify deleting depth three: within the
resolved boundary, depth three has real codimension two.

\begin{theorem}[Sharp codimension truncation]\label{P5a2}
\status{Proved.}
For every \(r\geq1\), the inclusion
\[
                         L_{\leq r}\longrightarrow L
\]
is \((r-1)\)-connected. Consequently depths at most \(k+2\), with the
full flag diagram, determine \(\tau_{\leq k}L\). The uniform bound is
optimal: for positive integral cutoffs \(r\), depth at most \(r\)
determines the \(k\)-type for every NC pair if and only if \(r\geq k+2\).

The assertion holds for analytic DM stacks on classifying spaces.
If every stratum is aspherical, the total normal-torus links in this
derived reconstruction can be replaced by classifying spaces of
their fundamental groupoids. This replacement does not require
meridian injectivity into the open-stratum group.
\end{theorem}

\begin{proof}
At a point of depth \(a\), a chart of \(X^{\log}\) has \(a\)
nonnegative radial coordinates and \(a\) circle coordinates. Within
its boundary the locus where those radial coordinates all vanish
has real codimension \(a-1\). The removed strata, of depths at least
\(r+1\), therefore have real codimension at least \(r\) in \(L\).

Triangulate the corner charts compatibly, or round their corners and
use the resulting locally flat strata. Relative general position moves
a map of a compact polyhedron of dimension less than \(r\) off the
removed union while fixing an already disjoint closed subpolyhedron.
Applied to spheres and nullhomotopies, this gives isomorphisms on
\(\pi_i\) for \(i<r-1\) and surjectivity on \(\pi_{r-1}\).
For noncompact manifolds it is applied to the compact image of each
map; no stable-end assertion is needed. For example, for
\(X=\C^2,D=\{z_1=0\}\), the normal boundary is
\(\C\times S^1\simeq S^1\), while the full end of
\(X^*=\C^*\times\C\cong S^1\times\R^3\) has cross-section
\(S^1\times S^2\). Thus even these homotopy types can differ.

For a stack choose an \'{e}tale atlas and its entire \v{C}ech nerve.
The manifold result applies to each nerve term. Geometric realization
preserves a fixed connectivity bound for maps of spaces. One proof
uses relative CW models: their relative cells have the required
minimum dimension, and realization adds a nonnegative simplicial
dimension. Equivalently, the appropriate left class of the
connected/truncated factorization system is closed under homotopy
colimits. Thus the termwise inclusions give the same bound on
classifying spaces. This argument includes nonproper analytic DM
stacks with infinite discrete isotropy, rather than assuming finite
orbifold charts.

For necessity take \(m=r+1\) coordinate hyperplanes in \(\C^m\).
The radial calculation above gives
\[
 L\simeq(S^1)^m,\qquad
 L_{\leq r}\simeq(S^1)^m\times S^{m-2}
                    =(S^1)^{r+1}\times S^{r-1}.
\]
For \(r=1\) this changes components; for \(r\geq2\) it changes
\(\pi_{r-1}\). Whenever \(r\leq k+1\), the \(k\)-types differ.
All these sharpness examples already have aspherical strata and
injective full meridian lattices. The supplied exact boundary-matrix
calculations verify the simplex boundary homology for dimensions
\(0,\ldots,4\); the radial proof establishes the assertion in every dimension.

Finally, a normal-torus bundle over an aspherical stratum is
aspherical by its homotopy exact sequence, as are the iterated flag
links. Objectwise replacement by \(B\Pi_1\) is consequently an
equivalence of coherent diagrams, and remains so under homotopy
colimits. This concerns a derived diagram of spaces; replacement of
the exit category by its ordinary homotopy category is a different
operation.
\end{proof}

\begin{corollary}[The corrected fundamental-group conclusion]\label{P5a3}
\status{Proved.}
Coherent fundamental-groupoid data from depths one, two, and three
always determine \(\Pi_1(L)\). Depths one and two suffice if there
are no triple crossings. If the collar inclusion
\(L\to X^*=X\setminus D\) is a \(1\)-type equivalence, the conclusion
is \(\Pi_1(X^*)\), and need not be \(\Pi_1(X)\).
\end{corollary}

\begin{proof}
Apply Theorem~\ref{P5a2} with \(k=1\). The fundamental-groupoid
functor is the left adjoint from spaces to one-types and therefore
preserves colimits, interpreted as bicategorical colimits in groupoids.
Consequently the full coherent diagram of fundamental groupoids always
computes \(\Pi_1(L)\), without an asphericity assumption on its pieces.
Without triple crossings, \(L_{\leq2}=L\). Compose with the assumed
equivalence to \(X^*\).
For the last distinction take \((X,D)=(\C,\{0\})\): its collar is
\(S^1\simeq\C^*\), but \(X\) is simply connected. Thus the unstarred
\(X\) visibly written in (247), scan 152, needs a mathematical
correction under this hypothesis; it is not silently transcribed
as \(X^*\).
\end{proof}

\subsection{Dimension filtration and its actual homotopy fibre}

Let \((g,n)\) be stable, \(d=3g-3+n\), and let \(\Gamma_{g,n}\)
fix the \(n\) labels individually. A multicurve means an isotopy
class of a finite family of pairwise disjoint, essential,
nonperipheral, mutually nonisotopic simple closed curves.
The empty multicurve is explicitly allowed. The stratum associated
to \(\sigma\) has dimension \(d-|\sigma|\). For every \(k\geq0\), put
\[
 P_k=\{\sigma:|\sigma|\geq d-k\},\qquad
 \sigma\leq\tau\ \Longleftrightarrow\ \sigma\subseteq\tau .
\]
Let \(W_k\) be the union, inside the open moduli stack, of deleted
tubes around strata of dimension at most \(k\). A precise model
requires the curves of some such \(\sigma\) to have hyperbolic
length less than a fixed sufficiently small \(\epsilon\).
For \(k\geq d\), the open stratum itself is retained, and
\(W_k=\mathcal M_{g,n}\).

\begin{theorem}[Fibre of the modular dimension filtration]\label{P5b1}
\status{Proved.}
For every stable \((g,n)\) and every \(k\geq0\), the homotopy fibre of
\[
                          BW_k\longrightarrow B\mathcal M_{g,n}
\]
is homotopy equivalent to \(|P_k|\). All stack stabilizers are retained.
For \(k\geq d\) the poset has the minimum element \(\varnothing\),
so its realization is contractible.
\end{theorem}

\begin{proof}
Work in marked Teichm\"uller space \(\mathcal T_{g,n}\).
Choose \(\epsilon\) so that intersecting essential geodesics cannot
both have length less than \(\epsilon\), as supplied by the collar
inequality. Define
\[
 V_\sigma=\{x\in\mathcal T_{g,n}:
               \ell_\gamma(x)<\epsilon\text{ for all }\gamma\in\sigma\}.
\]
Extend \(\sigma\) to a pants decomposition. Fenchel--Nielsen
coordinates express \(V_\sigma\) as a product of \((0,\epsilon)\)
and \((0,\infty)\) length intervals and real twist lines, so it is
contractible; see \citet[Chapter 1, Theorem 1.1]{Harer1988}.
The convention \(V_\varnothing=\mathcal T_{g,n}\) is essential.

Set \(V_k=\bigcup_{\sigma\in P_k}V_\sigma\). A finite intersection
of these open sets is empty for an incompatible union of curves,
and otherwise is \(V_{\cup\sigma}\); the union again belongs to
\(P_k\). This numerable good cover is closed under nonempty finite
intersections. Its homotopy-colimit construction gives an
equivariant span of underlying weak equivalences
\[
 V_k\ \longleftarrow\
 \hocolim_{\sigma\in P_k^{\mathrm{op}}}V_\sigma
 \ \longrightarrow\ |P_k^{\mathrm{op}}|.
\]
For the left arrow, the augmented \v{C}ech indexing nerve at a
point has a largest object: the full, finite system of its short
curves. It is therefore contractible. For the right arrow use the
contractibility of each \(V_\sigma\). Opposite posets have the same
order complex up to its canonical reversal.

The same length and twist coordinates describe a deleted pinching
tube: length gives the radial coordinate of the neck, twist its
angular coordinate, and the remaining marked component coordinates
its base. Sufficiently small bounds admit contractible refinements.
The span thus computes the specified tubular model of \(W_k\);
compare the boundary construction of
\citet[Chapter 3, Lemma 3.6]{Harer1988}.

Finally use the Borel constructions
\[
 B\mathcal M_{g,n}\simeq
 E\Gamma_{g,n}\times_{\Gamma_{g,n}}\mathcal T_{g,n}
 \simeq B\Gamma_{g,n},\qquad
 BW_k\simeq E\Gamma_{g,n}\times_{\Gamma_{g,n}}V_k .
\]
The second space has homotopy fibre \(V_k\) over \(B\Gamma_{g,n}\).
Contractibility of Teichm\"uller space proves the assertion for
the displayed map to \(B\mathcal M_{g,n}\).
A coarse quotient would discard stabilizer classifying spaces,
including ineffective inertia, and is not the construction used.
For \(k\geq d\), \(V_\varnothing\) already makes \(V_k\) the whole
Teichm\"uller space and verifies the endpoint directly.
\end{proof}

This poset remembers actual isotopy classes on a marked surface.
It is infinite and is not the finite poset of stable labelled
trees in a coarse boundary complex. Thus the finite calculations
for \(\Delta_{0,5}\) and \(\Delta_{0,6}\) do not alone verify this
homotopy fibre.

\begin{theorem}[Genus-zero connectivity at every level]\label{P5b2}
\status{Proved.}
For \(n\geq3\) and \(k\geq0\), the realization \(|P_k|\) for the
\(n\)-punctured sphere is \((k-1)\)-connected.
Consequently \(BW_k\to\mathcal M_{0,n}\) is \(k\)-connected.
For \(k\geq n-3\), the fibre is contractible.
\end{theorem}

\begin{proof}
The endpoint was proved above and includes \(n=3\).
Suppose \(0\leq k<d=n-3\). Let \(K=\mathcal C_{0,n}\) be the
unmodified disjointness curve complex. It is pure of dimension
\(d-1\), and the genus-zero case of
\citet[Chapter 4, Theorem 4.1]{Harer1988} identifies its homotopy
type as a wedge of \((d-1)\)-spheres.

For a face \(\sigma\) with \(q\) vertices, cutting along its curves
gives punctured spheres with \(n_v\) special points, and
\[
              \operatorname{Lk}_{K}(\sigma)
                       =\mathop{*}_{v}\mathcal C_{0,n_v}.
\]
The \(n_v=3\) factors are empty complexes, the join units \(S^{-1}\).
All other factors have the sphere dimensions just stated.
Since \(\sum_v(n_v-3)=d-q\), this link is
\((d-q-2)\)-connected, with homotopy dimension \(d-q-1\).
This establishes the homotopy Cohen--Macaulay property of all links.

Here is a rank-selection proof valid for the infinite complex.
Use the weak CW realization of the barycentric subdivision of \(K\).
Successively delete vertices corresponding to faces of sizes
\(1,2,\ldots,d-k-1\). When a face of size \(q\) is deleted, its
remaining link consists of proper cofaces and is the barycentric
subdivision of \(\operatorname{Lk}_{K}(\sigma)\). It is at least
\((k-1)\)-connected because \(q\leq d-k-1\).
The pre-deletion complex is obtained from the post-deletion complex
by coning this link. That inclusion is \(k\)-connected and
preserves homotopy through degree \(k-1\).
Vertices of the same rank are incomparable; their open stars are
disjoint, and the cones may be attached simultaneously.
This also works for infinitely many vertices: relative CW
connectivity is preserved under the coproduct of these cone
attachments, or one can apply the argument to the finite
subcomplex traversed by a compact homotopy in the weak CW topology.
The initial complex is \((d-2)\)-connected and the final one is
\(|P_k|\). For \(k=0\), its required nonemptiness follows directly
from a pants decomposition.

This is the homotopical rank-selection mechanism of
\citet[Corollary 4.4]{BjornerWachsWelker2009}, with the infinite
case justified explicitly instead of inferred from a finite-poset
statement.
\end{proof}

Here \(\mathcal C_{0,4}\) is an infinite discrete set; using a
modified Farey complex changes the question. The empty multicurve
also cannot be omitted. For example, at \((n,k)=(6,3)\) omission
would leave \(\mathcal C_{0,6}\), instead of the contractible fibre.

\begin{proposition}[Independent checks at \(n=5,6\)]\label{P5b3}
\status{Proved.}
The fibres for \((n,k)=(5,1)\) and \((6,2)\) are respectively
connected and simply connected by an independent argument on the
actual configuration spaces. The fibre for \((6,3)\) is contractible.
\end{proposition}

\begin{proof}
Normalize three points to \(0,1,\infty\). For \(d=n-3\),
\[
 \mathcal M_{0,n}
 =\operatorname{Conf}_{d}(\C\setminus\{0,1\})
 =\C^d\setminus\bigcup\{z_i=0,\ z_i=1,\ z_i-z_j=0\}.
\]
Let \(\mathcal L\) be these affine linear forms, and put
\[
 f(z)=\sum_i |z_i|^2+\sum_{\ell\in\mathcal L}|\ell(z)|^{-2}.
\]
The function is proper: the first term controls infinity and the
others diverge along every deleted hyperplane. If \(a_\ell\)
denotes the coefficient vector of \(\ell\), its Levi matrix is
\[
                  I_d+\sum_{\ell\in\mathcal L}
                            |\ell(z)|^{-4}a_\ell a_\ell^*,
\]
which is positive definite. A generic real linear perturbation
leaves this matrix unchanged and makes \(f\) Morse.
The critical points are the inverse image of a regular value
of its real gradient. That gradient is semialgebraic, so its
zero-dimensional regular inverse image is finite.
The perturbed function remains proper and bounded below.

At a critical point the real Hessian satisfies
\(Q(v)+Q(Jv)>0\) for nonzero \(v\). A negative subspace of
dimension larger than \(d\) would intersect its \(J\)-translate
nontrivially, contradicting that inequality.
Every Morse index is therefore at most \(d\).
A regular sublevel above all critical values is a compact
\(2d\)-manifold with handle indices at most \(d\). Reversing
the handles gives a relative decomposition from its boundary
with indices at least \(d\). Its boundary inclusion is consequently
\((d-1)\)-connected. For the precise handle statements see
\citet[Theorems 3.2 and 3.5, Remark 3.3, Theorem 7.2, and the
reversed-index proof of Theorem 7.4]{Milnor1963}; the exhaustion
and index estimate were supplied here.

Properness of \(\overline{\mathcal M}_{0,n}\) gives its NC
boundary a stable collar. The Morse level above all critical
values also has a stable collar. These are cofinal neighbourhoods
of the same end, and their product retractions identify their
homotopy types over \(\mathcal M_{0,n}\).
Thus this boundary computes \(W_{d-1}\) and its actual homotopy
fibre. For \(d=2\) that fibre is connected; for \(d=3\) it is
simply connected. At \(n=6,k=3\) the empty multicurve gives
contractibility by Theorem~\ref{P5b1}.

The exact check \texttt{levels\_stein.py} replaces
\(|\ell|^{-4}\) by independent positive variables and expands
every nonempty principal minor over \(\Z\).
For \(n=5\) it finds five forms and 14 positive determinant
monomials; for \(n=6\), nine forms and 93 positive determinant
monomials. Every minor has constant coefficient one and positive
remaining coefficients. This verifies the displayed exhaustion
algebra; the proof, rather than a finite quotient computation,
establishes connectivity of the infinite fibres.
\end{proof}

\begin{proposition}[The higher-genus extrapolation]\label{P5b4}
\status{Refuted.}
The assertion that \(|P_k|\) is always \((k-1)\)-connected in
every genus is false. A counterexample is \((g,n,k)=(2,1,3)\).
\end{proposition}

\begin{proof}
Here \(d=4\), so \(P_3\) contains all nonempty multicurves.
Its realization is the barycentric subdivision of
\(\mathcal C_{2,1}\). By
\citet[Chapter 4, Theorem 4.1 and Corollary 4.2]{Harer1988},
this complex has the homotopy type of a nontrivial wedge of
spheres of dimension \(2g-3+n=2\).
Its nonzero second homology is the dualizing module appearing
in that corollary. Hence it is not two-connected.
The all-genera fibre identification and the contractible
endpoint \(k\geq d\) remain valid. The all-genera two-level
statement \(k\leq2\), proved next, also remains valid.
\end{proof}

\subsection{Presentations and complete first identities}

\begin{theorem}[The two-dimensional presentation procedure in all genera]
\label{P5c1}
\status{Proved.}
For every stable \((g,n)\), \(P_0\) is nonempty, \(|P_1|\) is
connected, and \(|P_2|\) is simply connected.
The action of \(\Gamma_{g,n}\) on \(|P_2|\) therefore gives a
developable complex of groups with fundamental group
\(\Gamma_{g,n}\), retaining all stabilizers.
When \(d\geq3\) the quotient uses boundary strata of dimension
at most two and their incidence flags.
Stabilizers, incidence injections, transport elements, and
triangle relations give a presentation. Dimension-two
generators can be represented by lower-dimensional tube loops;
thus dimension one supplies generators and dimension two
supplies relations. The local groups are retained with their own
presentations at every level; the level-two assertion concerns the
additional gluing relations.
\end{theorem}

\begin{proof}
For \(k\geq d\) use the empty-multicurve cone. For \(k=0\)
a pants decomposition proves nonemptiness. It remains to
treat \(1\leq k\leq2<d\), or \(k=1<d\).
The sphere dimensions of \citet[Chapter 4, Theorem 4.1]{Harer1988}
show that a curve complex of complexity one is nonempty, one of
complexity two is connected, and one of complexity at least three
is simply connected. Explicitly the sphere dimensions are \(n-4\)
for \(g=0\), \(2g-3+n\) for \(g>0,n>0\), and \(2g-2\)
for closed surfaces; each has the asserted lower bound
in the indicated stable range.

The link of a multicurve is the join of the curve complexes
of its cut components. If the residual complexities are
\(d_i\), their sum is \(d-|\sigma|\); complexity-zero
components contribute join units. When that sum is at
least two, the join is connected: either one factor
already has complexity at least two, or two nonempty
factors occur. When it is at least three, the join is
simply connected: either one factor has complexity at
least three, a connected factor joins a nonempty factor,
or at least three nonempty factors occur. These assertions
follow from the join connectivity inequality.

Now delete vertices of the barycentric curve complex in
increasing face-size order, down to sizes at least \(d-k\).
At each deletion the remaining link is the upper link,
namely the curve complex of the cut surface just considered.
Its residual complexity is at least \(k+1\), so it is
\((k-1)\)-connected. Coning it is a \(k\)-connected
attachment. The initial curve complex already has this
connectivity, and the deletion argument of
Theorem~\ref{P5b2} applies. All realizations use the weak
CW topology; relative cone attachments, or finite compact
witnesses therein, justify the infinite family of deletions.
This proves the all-genera assertions at levels zero,
one, and two.

The action on the order complex has no inversions, because
the ranks of the vertices in a chain are distinct.
Its quotient has finitely many cells, represented by
stable labelled dual graphs and contraction flags.
Take a further barycentric subdivision and choose lifts
of its cells. Let \(G_v\) be the full stabilizer of each
chosen cell barycentre. A face incidence is oriented
from the larger cell to its face, giving after transport
an injection \(\psi_a:G_{i(a)}\to G_{t(a)}\).
For composable incidences choose the comparison element
\(g_{a,b}\in G_{t(a)}\), so that
\[
                  \psi_a\psi_b
                 =\operatorname{Ad}(g_{a,b})\psi_{ab}.
\]
Use generators of all \(G_v\)'s and symbols for the
incidence arrows, with the internal stabilizer relations,
\[
 ah a^{-1}=\psi_a(h),\qquad
 ab=g_{a,b}(ab)_{\mathrm{edge}},
\]
and set a maximal tree of edge symbols equal to one.
Reversed orientations give inverse symbols.
Every \(g_{a,b}\) is determined by comparing the specified
lifts; none is silently discarded.

For completeness, a path from a chosen lift to its
translate by \(\gamma\in\Gamma_{g,n}\) gives a word in
these transport and stabilizer generators. This proves
surjectivity. A word in the kernel lifts to a closed
edge path, which bounds a simplicial disc because
\(|P_2|\) is simply connected. Removing its triangles
uses the displayed triangle relations; changing lifts
at its corners uses the stabilizer and conjugation
relations. This proves injectivity, including for
ineffective elements which act trivially on the
underlying complex but remain in its local groups.

The stabilizers are the mapping class groups of cut
surfaces with gluing twists retained, together with
the finite symmetries of their dual graphs.
The subgroup preserving all components and branches
is an extension of the product of component mapping
class groups by the gluing-twist lattice.
This extension and its coherence are retained; a
general-genus splitting is not asserted.
The component dimensions sum to at most two, so the
only positive-dimensional indecomposable factors
are of types \((0,4),(1,1),(0,5),(1,2)\).
At a dimension-two factor the already proved
connectedness of its \(P_1\) means that its boundary
tube maps onto its fundamental group. For a product
of two dimension-one factors, the join of their
nonempty boundary fibres is connected, giving the
same surjectivity. This remains true equivariantly
under graph symmetries and after taking normal-torus
bundles. Thus dimension-two generators can be replaced
by lower-dimensional words, leaving their relations.

In genus zero the dual graphs are trees, pure labels
remove the graph permutations, and framings split the
twist extensions noncanonically. The positive-dimensional
factors are then \(\mathcal M_{0,4}\),
\(\mathcal M_{0,5}\), or \(\mathcal M_{0,4}^2\).
For \(n=5\), \(P_2\) contains the empty vertex with
stabilizer \(\Gamma_{0,5}\); the abstract presentation
would be tautological there. The following independent
calculation supplies that elementary group.
\end{proof}

Write \(F_m\) for a free group of rank \(m\), and use
\(c(u,v)=v^{-1}u^{-1}vu\) throughout.
For ordinary braid generators define
\[
 A_{ij}=\sigma_{j-1}\cdots\sigma_{i+1}\,
             \sigma_i^2\,\sigma_{i+1}^{-1}
                         \cdots\sigma_{j-1}^{-1}
                  \quad (i<j).
\]
The first two finite marked points are fixed at
\(0,1\); the last spherical marked point is \(\infty\).
The remaining points move in that punctured affine line.

\begin{proposition}[An explicit presentation of \(\Gamma_{0,5}\)]
\label{P5c2}
\status{Proved.}
One has
\[
                \Gamma_{0,5}=F(x,y,z)\rtimes F(a,b),
\]
where \(a=A_{13}, b=A_{23}, x=A_{14}, y=A_{24}, z=A_{34}\).
A complete presentation consists of these five generators
and the six relations
\[
\begin{array}{lll}
 axa^{-1}=z^{-1}xz,&
 aya^{-1}=c(x,z)y c(x,z)^{-1},&
 aza^{-1}=c(x,z)z,\\[2pt]
 bxb^{-1}=x,&
 byb^{-1}=z^{-1}yz,&
 bzb^{-1}=c(y,z)z .
\end{array}
\]
The presentation complex is aspherical, with cell
counts \((1,5,6)\).
\end{proposition}

\begin{proof}
Forget the second moving point of
\(\operatorname{Conf}_2(\C\setminus\{0,1\})\).
The base is \(\C\setminus\{0,1\}\), with free
generators \(a,b\), and the fibre is
\(\C\setminus\{0,1,z_1\}\), with free generators
\(x,y,z\). This is the locally trivial fibration of
\citet[Theorem 1]{FadellNeuwirth1962}.
A section places the last point on the positive
real axis farther from the origin than every
preceding point and than \(1\). Base and fibre
are aspherical, so the homotopy exact sequence
gives the split group extension and asphericity
of its total space.

Our Artin convention is
\(\sigma_i(X_i)=X_iX_{i+1}X_i^{-1}\),
\(\sigma_i(X_{i+1})=X_i\), with other letters fixed.
For the stated kernel generators, conjugation
by a positive base braid is the Artin action of
its mirror word, obtained by negating each
Artin-generator exponent. Substituting \(A_{13}\)
and \(A_{23}\) gives the six displayed formulae.
The program \texttt{levels\_braids.py} compares
each equality using the actual braid words in
the free-group automorphism representation on
four letters. That representation is faithful;
see \citet[Theorem A, representation (1),
parameter \(1\)]{Shpilrain2001}.
The checks concern full noncommutative words.

Replace the base by a two-petalled rose and
the fibre by a three-petalled rose. Realize the
two monodromies by based graph homotopy
equivalences and form their mapping cylinders
over the base loops. The resulting graph bundle
up to homotopy has one vertex, five edges,
and exactly the six indicated two-cells.
Its universal cover is a tree of contractible
fibre trees and is contractible. The displayed
presentation is therefore complete and
aspherical.
\end{proof}

\begin{proposition}[A presentation and all first syzygies for
\(\Gamma_{0,6}\)]\label{P5c3}
\status{Proved.}
Adjoin \(p=A_{15},q=A_{25},r=A_{35},s=A_{45}\).
Then
\(\Gamma_{0,6}=F(p,q,r,s)\rtimes\Gamma_{0,5}\).
Keep the six base relators and add the twenty
relations \(huh^{-1}=\phi_h(u)\) in this table:
\[
\renewcommand{\arraystretch}{1.5}
\begin{array}{c|cccc}
h&\phi_h(p)&\phi_h(q)&\phi_h(r)&\phi_h(s)\\ \hline
a&r^{-1}pr&c(p,r)q c(p,r)^{-1}&c(p,r)r&s\\
b&p&r^{-1}qr&c(q,r)r&s\\
x&s^{-1}ps&c(p,s)q c(p,s)^{-1}
                 &c(p,s)r c(p,s)^{-1}&c(p,s)s\\
y&p&s^{-1}qs&c(q,s)r c(q,s)^{-1}&c(q,s)s\\
z&p&q&s^{-1}rs&c(r,s)s
\end{array}
\]
This gives nine generators and twenty-six
relators and a finite aspherical three-complex
with cell counts \((1,9,26,24)\).
The twenty-four identities specified below,
one for each of the six base relators and
each of \(p,q,r,s\), freely generate all
first identities modulo Peiffer identities,
equivalently the \(\pi_2\) of this presentation
complex, over \(\Z\Gamma_{0,6}\).
\end{proposition}

\begin{proof}
Forget the third moving point. The fibre is
the affine line with four points removed,
and the far-out-point section splits the
Fadell--Neuwirth extension. Substitution in
the mirrored Artin action gives the table.
Every entry is checked in
\texttt{levels\_braids.py} against actual
conjugation of \(A_{i5}\) in the five-strand
braid group. Each of the six base relators
also acts trivially on all four fibre
generators, giving twenty-four exact
free-word checks. Thus the table defines
the required action of \(\Gamma_{0,5}\);
the split extension proves completeness
of the twenty-six defining relations.

Here are the full identity words as a
finite recursive schema. Temporarily
regard all nine generators as free.
For a base word \(w\) and a fibre word
\(v\), put
\[
                  D_w(v)=wvw^{-1}\phi_w(v)^{-1}.
\]
For positive generators \(h,u\), \(D_h(u)\)
is its mixed relator. Expand other
\(D\)-words as products of conjugates
of those relators using the exact rules
\[
\begin{split}
 D_{ab}(v)&=aD_b(v)a^{-1}D_a(\phi_b(v)),\\
 D_h(vv')&=D_h(v)\phi_h(v)D_h(v')\phi_h(v)^{-1},\\
 D_h(v^{-1})&=\phi_h(v)^{-1}D_h(v)^{-1}\phi_h(v),\\
 D_{h^{-1}}(v)&=h^{-1}
                    D_h(\phi_h^{-1}(v))^{-1}h .
\end{split}
\]
The inverse automorphisms are fixed by the
mirrored-braid convention and supplied as
integer words in
\texttt{levels\_braids\_actions.json}.
Let \(R\) run through the six base words
\(hvh^{-1}(\text{its displayed image})^{-1}\),
with \(h=a,b\) and \(v=x,y,z\).
For each new generator \(u=p,q,r,s\),
\(\phi_R(u)=u\) freely, and the identity is
\[
                   R\,(uR^{-1}u^{-1})\,D_R(u)^{-1}=1. \tag{*}
\]
The first two factors are a base relator
and a conjugate inverse base relator;
the last expands into conjugates of the
twenty mixed relators by the stated rules.
This specifies every word and sign.
The program \texttt{levels\_syzygies.js} records each of
these twenty-four products and verifies
that it freely reduces to the empty word.
The products have between eight and 640
conjugate factors; the largest conjugating
word has length 171.

A check that identities hold does not
establish their completeness. For that,
let \(K_5\) be the aspherical two-complex
just constructed, and \(R_4\) a
four-petalled rose. Each component of
\[
              \Map_*(R_4,R_4)\simeq(\Omega R_4)^4
\]
is contractible, since \(R_4\) is
aspherical. Based self-homotopy
equivalences form the components indexed
by \(\operatorname{Aut}(F_4)\).
The exact action therefore constructs
a graph bundle up to homotopy over
\(K_5\), with a section. Choose the
graph maps over base edges and the
relator homotopies over base two-cells.
Contractibility of the mapping-space
components supplies their coherence.

Over a base \(q\)-cell attach its
section \(q\)-cell first. The remaining
relative pair is
\[
 (D^q\times R_4,\
   S^{q-1}\times R_4\ \cup\ D^q\times\{*\}),
\]
with quotient \(S^q\wedge R_4\), a
wedge of four \((q+1)\)-spheres.
Thus the base counts \((1,5,6)\) give
\[
 c_0=1,\quad c_1=5+4=9,\quad
 c_2=6+4\cdot5=26,\quad c_3=4\cdot6=24.
\]
The additional two-cells are the
twenty mixed relators. A base relator
disc crossed with a fibre edge has
boundary the sphere \((*)\).
The total space is aspherical by the
homotopy exact sequence of this
homotopy bundle: its base and fibre
have no higher homotopy, and its
group is the specified extension.

Let \(G=\Gamma_{0,6}\) and pass to
universal covers. There are no
cells above degree three, and
contractibility of the full cover
gives an exact sequence
\[
 0\longrightarrow \Z G^{24}
       \xrightarrow{\partial_3}\Z G^{26}
       \xrightarrow{\partial_2}\Z G^9 .
\]
Hence \(\partial_3\) is injective
onto \(\ker\partial_2\).
The universal cover of the
presentation two-skeleton is
simply connected, so Hurewicz
identifies its \(\pi_2\) with
that kernel. The twenty-four
attaching spheres are consequently
a free \(\Z G\)-basis of all
first identities modulo Peiffer
identities.
\end{proof}

This last module belongs to the
displayed finite presentation.
It is not identified with
\(\pi_2\) of the full geometric
boundary tube or of the infinite
curve complex, which have different
cell models.

\begin{theorem}[Two levels, their elementary cases, and dimension-three
identities]\label{P5c4}
\status{Proved.}
In every genus, strata of modular
dimension at most two suffice
for the ordinary fundamental-group
presentation, with full incidence
and stabilizer data.
The four indecomposable positive
dimensional types in folio 29,
scan 60, are
\[
                         (0,4),\ (1,1),\ (0,5),\ (1,2).
\]
The first two have dimension one,
the last two dimension two.

In genus zero, dimension-three
coherence has the support types
\[
 \mathcal M_{0,6},\qquad
 \mathcal M_{0,5}\times\mathcal M_{0,4},\qquad
 \mathcal M_{0,4}^3,
\]
with zero-dimensional factors allowed.
The finite presentations specified
above and below have respectively
24, 12, and 8 complete first identities.
\end{theorem}

\begin{proof}
The all-genera group presentation
is Theorem~\ref{P5c1}.
Solving \(3g-3+n=1\) in the stable
range gives \((0,4),(1,1)\);
solving \(3g-3+n=2\) gives
\((0,5),(1,2)\).
For an independent published
topological formulation,
\citet[Section 2, Theorem 2,
and Remark 4]{HatcherLochakSchneps2000}
prove simple connectivity of
the maximal-curve-system complex
with elementary moves and
relations and identify these
four support types. Their
disjoint-support commutation
relations are included.
This is an ordinary topological
two-level statement, not a
profinite or Galois reconstruction
assertion.

In genus zero each stratum is
a product of \(\mathcal M_{0,n_v}\),
with dimensions adding.
The partitions of three give
exactly the three listed types.
Theorem~\ref{P5b2} at \(k=3\)
says that this level kills
the remaining \(\pi_2\)
obstructions to its comparison.
Proposition~\ref{P5c3} gives
the first type explicitly.

For the second type use the
five-generator presentation
of \(\Gamma_{0,5}\) and add
two generators \(t_1,t_2\).
Keep its six relators and
impose the ten commutators
\([h,t_i]=1\), where
\(h=a,b,x,y,z\).
The aspherical model is
\(K_5\times R_2\), with
cell counts
\[
                              (1,7,16,12).
\]
Its twelve first identities
are \((*)\) for the six
base relators and the two
new letters, with trivial
action \(\phi_h(t_i)=t_i\).

For the third type take
generators \(a_1,a_2\),
\(b_1,b_2\), \(c_1,c_2\)
and all twelve commutators
between distinct pairs of
these three free factors.
The model \(R_2^3\) has
cell counts
\[
                               (1,6,12,8).
\]
Its eight cube identities
are the same schema for
the four base relators
\(a_i b_j a_i^{-1}b_j^{-1}\)
and each \(c_\ell\).
Both product models are
aspherical and have no cells
above dimension three.
The universal-cover and
Hurewicz argument of
Proposition~\ref{P5c3}
therefore proves that the
12 and 8 identities are
complete free bases for
their presentation modules.
The expanded-word check
verifies all of them;
the second family has
six to fourteen conjugate
factors per identity,
and the cube identities
have six each.

Normal tori and lower
pieces can already have
internal higher cells.
Thus this support statement
does not say that every
three-cell first appears at
modular dimension three.
Nor does it extend the
higher-connectivity assertion
to every genus: the
counterexample in
Proposition~\ref{P5b4}
explains the precise limit.
\end{proof}

\paragraph{Punctures and fixed boundary circles.}
The explicit sphere groups fix
punctures individually. If instead
each of the \(n\) boundary circles
is fixed pointwise, capping
introduces central boundary
twists. For \(g=0,n\geq3\)
the result is noncanonically
\(\Gamma_{0,n}\times\Z^n\).
Indeed normalization at
\(0,1,\infty\) trivializes
all individual tangent-direction
circle bundles: use \(dz\)
at finite marks and \(d(1/z)\)
at infinity. Their product
is the framed moduli space,
whose fundamental group is
the pointwise-boundary mapping
class group. Different framings
alter the splitting.

\paragraph{Computational scope.}
The supplied exact programs
verify the nonabelian hexagon
presentation, twenty-six braid
equalities, twenty-four action
compatibilities, all 44
expanded first identities for
the three dimension-three
presentations, and the Levi
principal minors.
They do not infer infinite
curve-complex connectivity
from a finite stable-tree
quotient. The connectivity
and the completeness of the
syzygies are established by
the proofs above.
